\documentclass[11pt]{article}

\usepackage[margin=1in]{geometry}
\usepackage{amsmath, amssymb, amsthm}
\usepackage{booktabs}
\usepackage{mathtools}
\usepackage{hyperref}

\newtheorem{remark}{Remark}
\newcommand{\R}{\mathbb{R}}
\newcommand{\Om}{\Omega}
\newcommand{\eps}{\varepsilon}
\newcommand{\kk}{\kappa}
\newcommand{\ustar}{u^\star}
\DeclareMathOperator{\diag}{diag}
\DeclareMathOperator{\spanop}{span}

\title{From residual least-squares to supervised least-squares:\\
a family of preconditioners for the discrete Poisson problem}
\author{Marius Zeinhofer}
\date{\today}

\begin{document}
\maketitle

\begin{abstract}
We study a one-parameter family of preconditioners $P$ that interpolates between the
identity $P=I$ and the exact inverse $P=A^{-1}$ of the discrete Poisson stiffness matrix
$A$. At one end the preconditioned problem is the raw residual least-squares problem
$\tfrac12\|Au-b\|^2$; at the other it is the ``supervised'' least-squares problem
$\tfrac12\|u-\ustar\|^2$. We recall the exact eigenstructure of $A$ on a regular grid,
clarify \emph{which} part of the spectrum drives the ill-conditioning, and analyse two
natural interpolations---a \emph{convex blend} $P_t=(1-t)I+tA^{-1}$ and a
\emph{fractional power} $P_s=A^{-s}$---comparing how each reshapes the spectrum, their
conditioning under mesh refinement, and their interpretation as a descent direction.
No neural networks appear here; everything is a statement about the conditioning of a
linear least-squares problem.
\end{abstract}

\section{The discrete Poisson operator}

Let $\Om=(0,1)^2$ and consider the model problem
\begin{equation}
  -\Delta u = f \quad\text{in }\Om, \qquad u=0 \quad\text{on }\partial\Om .
\end{equation}
Let $V_h\subset H_0^1(\Om)$ be the space of continuous piecewise-bilinear ($Q_1$) finite
elements on a \emph{regular} grid, with nodal basis $\{\varphi_i\}_{i=1}^{N}$ on the
interior nodes. The (symmetric positive definite) \emph{stiffness matrix} is
\begin{equation}
  A_{ij} \;=\; \int_\Om \nabla\varphi_i \cdot \nabla\varphi_j \, dx ,
  \qquad i,j = 1,\dots,N ,
  \label{eq:stiffness}
\end{equation}
and the discrete solution $\ustar\in\R^{N}$ (the vector of nodal values) solves
$A\,\ustar = b$ for a load vector $b$. We do not need the precise form of $b$ here; the
analysis concerns the conditioning of $A$ and of the preconditioned problems built from it.

Throughout, $A$ is SPD, so it admits an orthonormal eigendecomposition
\begin{equation}
  A = Q\Lambda Q^\top,\qquad \Lambda = \diag(\lambda_1,\dots,\lambda_N),\qquad
  0<\lambda_{\min}\le\cdots\le\lambda_{\max}.
\end{equation}
Every operator we consider is a function of $A$ (hence shares the eigenvectors $Q$), so we
can track it eigenvalue-by-eigenvalue.

\section{Exact eigenstructure on a regular grid}

The regularity of the grid is, for the story below, \emph{secondary}: it only serves to
give closed-form eigenvalues that make the qualitative statements transparent. On a general
mesh the same statements hold with $\lambda_{\min},\lambda_{\max}$ replaced by their
numerical values.

\paragraph{One dimension.}
On $[0,1]$ with $n$ interior nodes and spacing $h=1/(n+1)$, the $Q_1$ (here: linear)
stiffness and mass matrices are the symmetric tridiagonal Toeplitz matrices
\begin{equation}
  K = \tfrac{1}{h}\,\mathrm{tridiag}(-1,\,2,\,-1),
  \qquad
  M = \tfrac{h}{6}\,\mathrm{tridiag}(1,\,4,\,1).
\end{equation}
Both are diagonalised by the \emph{same} discrete sine modes
$v^{(k)}_j = \sin\!\big(k\pi j h\big)$, $j,k=1,\dots,n$. Writing
$\theta_k = k\pi h = \dfrac{k\pi}{n+1}$, their eigenvalues are
\begin{align}
  \mu_k \;&=\; \frac{1}{h}\big(2 - 2\cos\theta_k\big)
            \;=\; \frac{4}{h}\,\sin^2\!\Big(\frac{\theta_k}{2}\Big)
            \qquad&&\text{(stiffness)},\\[2pt]
  \nu_k \;&=\; \frac{h}{6}\big(4 + 2\cos\theta_k\big)
            \;=\; \frac{h}{3}\big(2 + \cos\theta_k\big)
            \qquad&&\text{(mass)}.
\end{align}

\paragraph{Two dimensions.}
On the tensor-product $Q_1$ space the Laplacian stiffness separates as
\begin{equation}
  A \;=\; K\otimes M \;+\; M\otimes K ,
\end{equation}
which is exactly the assembled $Q_1$ ``nine-point'' Laplacian (its centre stencil value is
$8/3$). Its eigenvectors are the tensor sine modes $v^{(k)}\otimes v^{(l)}$, and its
eigenvalues are
\begin{equation}
  \boxed{\;
  \Lambda_{k,l} \;=\; \mu_k\,\nu_l \;+\; \nu_k\,\mu_l ,
  \qquad k,l=1,\dots,n .
  \;}
  \label{eq:eig2d}
\end{equation}

\paragraph{Asymptotics: which end is the culprit?}
Expanding \eqref{eq:eig2d} at the extreme modes (with $h\to0$):
\begin{align}
  \text{lowest mode }(k,l)=(1,1):\quad
    &\Lambda_{1,1} \approx 2\,\mu_1\nu_1 \approx 2\big(\pi^2 h\big)(h) = 2\pi^2 h^2 = O(h^2),\\
  \text{highest mode }(k,l)=(n,n):\quad
    &\Lambda_{n,n} \approx 2\,\mu_n\nu_n \approx 2\Big(\tfrac{4}{h}\Big)\Big(\tfrac{h}{3}\Big)
       = \tfrac{8}{3} = O(1).
\end{align}
Hence for the \emph{assembled stiffness matrix} \eqref{eq:stiffness},
\begin{equation}
  \lambda_{\min} = O(h^2)\;\longrightarrow\;0,
  \qquad
  \lambda_{\max} = O(1),
  \qquad
  \kk(A) = \frac{\lambda_{\max}}{\lambda_{\min}} = O\!\big(h^{-2}\big).
  \label{eq:condA}
\end{equation}
So it is the \textbf{small eigenvalues that vanish}, and they belong to the
\emph{low-frequency, smooth} modes $v^{(1)}\otimes v^{(1)}$, etc. The largest eigenvalue
stays $O(1)$.

\begin{remark}[Reconciling with the ``$-\Delta$ has large high-frequency eigenvalues'' intuition]
The continuous operator $-\Delta$ is unbounded \emph{above}: its eigenvalues
$\pi^2(k^2+l^2)$ grow with frequency, so one naturally thinks ``high frequency $=$ large
eigenvalue.'' That intuition refers to the \emph{generalised} eigenproblem
$A v = \lambda\,M v$, whose eigenvalues $\lambda_{k,l}=\Lambda_{k,l}/\text{(mass scaling)}
\approx \pi^2(k^2+l^2)$ do grow with frequency and are the physical (mesh-independent)
Rayleigh quotients. The \emph{matrix} conditioning that governs iterative solvers and
gradient descent, however, is set by the plain eigenvalues $\Lambda_{k,l}$ of $A$ in
\eqref{eq:eig2d}. There the overall $O(1)$ normalisation of the $Q_1$ stiffness places the
high-frequency modes near $O(1)$ and lets the low-frequency modes shrink to $O(h^2)$. The
condition number $\kk(A)=O(h^{-2})$ is the same regardless of normalisation---scaling $A$
by a constant (e.g.\ the $1/h^2$ of a finite-difference Laplacian would push
$\lambda_{\max}$ to $O(h^{-2})$ instead) never changes a ratio. What is invariant, and what
matters, is:
\begin{center}
\emph{the ill-conditioning is driven by the smooth, low-frequency modes, whose eigenvalues
are smallest.}
\end{center}
These are precisely the slowly-converging directions of gradient descent and the modes that
multigrid's coarse-grid correction is designed to handle.
\end{remark}

\section{Residual versus supervised least squares}

Given any nonsingular $P$ (a \emph{preconditioner}, i.e.\ an approximation of $A^{-1}$),
define the preconditioned residual least-squares functional
\begin{equation}
  J_P(u) \;=\; \tfrac12\,\big\| P\,(Au-b) \big\|^2
          \;=\; \tfrac12\,(Au-b)^\top P^\top P\,(Au-b).
  \label{eq:JP}
\end{equation}
Because $P$ is nonsingular, every $J_P$ has the \emph{same} minimiser $u=\ustar$; only the
\emph{geometry} differs. Its gradient and Hessian are
\begin{equation}
  \nabla J_P(u) = A P^\top P\,(Au-b) = H_P\,(u-\ustar),
  \qquad
  H_P \;=\; A\,P^\top P\,A .
  \label{eq:hessian}
\end{equation}
Two limiting cases bracket the family:
\begin{itemize}
  \item $P=I$: $\;J_I(u)=\tfrac12\|Au-b\|^2$ is the \textbf{residual least-squares}
        problem, with $H_I = A^2$ and $\kk(H_I)=\kk(A)^2 = O(h^{-4})$.
  \item $P=A^{-1}$: since $Au-b = A(u-\ustar)$,
        \begin{equation}
          J_{A^{-1}}(u) = \tfrac12\,(Au-b)^\top A^{-2}(Au-b)
                        = \tfrac12\,\|u-\ustar\|^2 ,
        \end{equation}
        the \textbf{supervised least-squares} problem, with $H_{A^{-1}}=I$ and $\kk=1$.
\end{itemize}
Thus $P$ tunes the problem continuously from ``match the residual'' to ``match the
solution.'' Because $P$ commutes with $A$, if $P$ has eigenvalue $p(\lambda)$ on the mode of
$A$ with eigenvalue $\lambda$, then the \emph{preconditioned operator} $PA$ has eigenvalue
\begin{equation}
  g(\lambda) \;=\; \lambda\,p(\lambda),
\end{equation}
and $H_P=A P^2 A$ has eigenvalue $g(\lambda)^2$. Consequently
$\kk(H_P)=\kk(PA)^2$, and it suffices to understand the spectrum of $PA$, i.e.\ the map
$\lambda\mapsto g(\lambda)$: perfect preconditioning is $g\equiv1$ (the flat spectrum), no
preconditioning is $g(\lambda)=\lambda$.

It is also useful to read $P$ as a \emph{descent direction}. The preconditioned
(Richardson) step on the residual is $u\mapsto u-\eta\,P(Au-b)$, and since
$A^{-1}(Au-b)=u-\ustar$,
\begin{equation}
  P\,(Au-b) \;=\; \underbrace{(Au-b)}_{\text{residual direction}}
      \ \text{when }P=I,
  \qquad
  P\,(Au-b) \;=\; \underbrace{(u-\ustar)}_{\text{exact error / ``supervised'' direction}}
      \ \text{when }P=A^{-1}.
  \label{eq:directions}
\end{equation}

\section{Two interpolations between $I$ and $A^{-1}$}

We compare two one-parameter families, both SPD, both functions of $A$, both with endpoints
$P=I$ and $P=A^{-1}$.

\subsection{Convex blend}
\begin{equation}
  \boxed{\,P_t \;=\; (1-t)\,I \;+\; t\,A^{-1}, \qquad t\in[0,1].\,}
\end{equation}
Its eigenvalue on the mode $\lambda$ is $p_t(\lambda)=(1-t)+t/\lambda$, so
\begin{equation}
  g_t(\lambda) = \lambda\,p_t(\lambda) = (1-t)\,\lambda + t,
  \qquad
  \text{$H_{P_t}$ eigenvalue } = \big((1-t)\lambda+t\big)^2 .
\end{equation}
The preconditioned spectrum is an \textbf{affine map} of $\{\lambda_i\}$: scale by $(1-t)$
and \emph{shift up by $t$}. The shift acts as an additive floor that lifts the dangerous
small eigenvalues ($\lambda_{\min}=O(h^2)$) up to $\approx t$. Its condition number is
\begin{equation}
  \kk(P_tA) = \frac{(1-t)\lambda_{\max}+t}{(1-t)\lambda_{\min}+t}
  \;\xrightarrow[h\to0]{}\; \frac{(1-t)\lambda_{\max}+t}{t} = O(1/t),
\end{equation}
\emph{bounded, hence mesh-independent, for every fixed $t>0$}. The blend therefore behaves
like Tikhonov regularisation: any $t>0$ already removes the $h$-blow-up, and increasing $t$
lowers the (bounded) condition number toward $1$.

Its descent direction is a \textbf{convex combination of the two endpoint directions} in
\eqref{eq:directions}:
\begin{equation}
  P_t\,(Au-b) \;=\; (1-t)\,(Au-b) \;+\; t\,(u-\ustar).
  \label{eq:blenddir}
\end{equation}

\subsection{Fractional power}
\begin{equation}
  \boxed{\,P_s \;=\; A^{-s}, \qquad s\in[0,1].\,}
\end{equation}
Its eigenvalue is $p_s(\lambda)=\lambda^{-s}$, so
\begin{equation}
  g_s(\lambda) = \lambda^{1-s},
  \qquad
  \text{$H_{P_s}$ eigenvalue } = \lambda^{2(1-s)} .
\end{equation}
This is a \textbf{uniform contraction of the log-spectrum}: $\log g_s(\lambda) =
(1-s)\log\lambda$, so every mode's logarithmic distance from $1$ is scaled by the same
factor $(1-s)$. The condition number is
\begin{equation}
  \kk(P_sA) = \kk(A)^{\,1-s} = O\!\big(h^{-2(1-s)}\big),
  \qquad
  \kk(H_{P_s}) = \kk(A)^{2(1-s)} .
\end{equation}
For every $s<1$ this still \emph{grows} under refinement (only the exponent is softened);
mesh-independence is reached only at the endpoint $s=1$. The parameter $s$ is a
dimensionless ``fraction of the log-conditioning removed'': at $s$ one has eliminated a
fraction $s$ of the exponent, uniformly across resolutions.

Its descent direction is a \textbf{spectrally reweighted error}:
\begin{equation}
  P_s\,(Au-b) \;=\; A^{-s}\,A\,(u-\ustar) \;=\; A^{1-s}\,(u-\ustar),
\end{equation}
i.e.\ each eigenmode of the exact error $u-\ustar$ is scaled by $\lambda^{1-s}$.

\subsection{Comparison}

\begin{center}
\renewcommand{\arraystretch}{1.35}
\begin{tabular}{@{}lll@{}}
\toprule
 & \textbf{Convex blend } $P_t=(1-t)I+tA^{-1}$ & \textbf{Power } $P_s=A^{-s}$ \\
\midrule
Preconditioned eig.\ $g(\lambda)$ & $(1-t)\lambda + t$ (affine) & $\lambda^{1-s}$ (power) \\
Effect on spectrum & additive floor on small $\lambda$ & uniform log-contraction \\
$\kk(PA)$ & $\dfrac{(1-t)\lambda_{\max}+t}{(1-t)\lambda_{\min}+t}$ & $\kk(A)^{1-s}$ \\
Mesh behaviour ($h\to0$) & bounded $O(1/t)$ for any $t>0$ & $O(h^{-2(1-s)})$, blows up if $s<1$ \\
Descent direction & $(1-t)(Au-b)+t(u-\ustar)$ & $A^{1-s}(u-\ustar)$ \\
Reading & convex mix of the two schemes & spectral filtering of the error \\
Endpoints & $t=0:\,I$;\quad $t=1:\,A^{-1}$ & $s=0:\,I$;\quad $s=1:\,A^{-1}$ \\
\bottomrule
\end{tabular}
\end{center}

\paragraph{Summary of the mathematical distinction.}
Both families connect the residual least-squares problem ($P=I$) to the supervised
least-squares problem ($P=A^{-1}$), and both flatten the preconditioned spectrum
$g(\lambda)$ from $\lambda$ to $1$. They differ in \emph{how}:
\begin{itemize}
  \item the \textbf{convex blend} reshapes the spectrum \emph{affinely}, flooring the small
        (smooth-mode) eigenvalues; any $t>0$ already yields a bounded, mesh-independent
        condition number, and its descent direction is the literal convex combination
        \eqref{eq:blenddir} of the residual and the exact-error directions;
  \item the \textbf{fractional power} reshapes the spectrum by a \emph{uniform log-contraction},
        giving the clean closed form $\kk(PA)=\kk(A)^{1-s}$; mesh-independence is attained
        only in the limit $s=1$, so across resolutions it produces a fan of conditioning
        slopes $2(1-s)$ (for $\kk(H_P)$) that interpolates between $O(h^{-4})$ at $s=0$ and
        $O(1)$ at $s=1$.
\end{itemize}
The blend is the more direct embodiment of ``interpolating between two training schemes'';
the power is the cleaner spectral dial and carries the sharper statement about behaviour
under mesh refinement.

\section{Numerical realization}
\label{sec:numerical}

Both families are functions of $A$, so both are realized by a \emph{single}
eigendecomposition, which is what makes them cheap to implement and to compare.

\paragraph{Interior operator.}
The assembled finite-element stiffness carries rows and columns for the Dirichlet boundary
nodes, but those nodes are not degrees of freedom (they are fixed to zero). The object of
interest is therefore the symmetric positive-definite \emph{interior block}, which is
precisely the matrix $A\in\R^{N\times N}$ analysed above ($N$ = number of interior nodes).
In practice one restricts the assembled matrix to the interior indices and works with that
block.

\begin{remark}[Interior space versus full-grid storage]
In the analysis $r\in\R^N$ has \emph{only} interior components: it is the coordinate vector
of the residual functional $v\mapsto \langle f,v\rangle - a(u_h,v)\in V_h'$ tested against
the interior basis functions, and it has no boundary component at all---the boundary carries
no test function in $H_0^1$. A finite-element or operator-learning implementation, however,
typically stores every field on the \emph{full} grid, boundary nodes included. There,
interior quantities are represented by the trivial extension by zero
$\R^{N}\hookrightarrow\R^{N_{\mathrm{full}}}$, $r\mapsto(r,0)$, and the assembled matrix is
restricted to its interior block. The boundary zeros are then pure storage
bookkeeping for this embedding; they do \emph{not} express that the residual ``vanishes on
the boundary''. All statements below are made in the interior space $\R^N$.
\end{remark}

\paragraph{Spectral construction.}
Compute once the symmetric eigendecomposition
\begin{equation}
  A = Q\Lambda Q^\top, \qquad Q^\top Q = I, \qquad \Lambda=\diag(\lambda_1,\dots,\lambda_N).
\end{equation}
For a scalar \emph{spectral map} $p:(0,\infty)\to\R$ applied entrywise to the eigenvalues,
define the preconditioner
\begin{equation}
  P \;=\; Q\,p(\Lambda)\,Q^\top,
  \qquad p(\Lambda) := \diag\big(p(\lambda_1),\dots,p(\lambda_N)\big).
\end{equation}
The two families are the two choices of $p$:
\begin{equation}
  \text{convex blend:}\quad p_t(\lambda) = (1-t) + \frac{t}{\lambda},
  \qquad\qquad
  \text{fractional power:}\quad p_s(\lambda) = \lambda^{-s}.
\end{equation}
The endpoints are exact and require no special-casing: $p\equiv 1$ gives $P=I$, and
$p(\lambda)=1/\lambda$ gives $P=Q\Lambda^{-1}Q^\top=A^{-1}$.

\paragraph{Application to a residual.}
Given a residual $r\in\R^{N}$, the action $Pr$ is
\begin{equation}
  Pr \;=\; Q\Big( p(\Lambda)\odot \big(Q^\top r\big)\Big),
  \label{eq:apply}
\end{equation}
i.e.\ (i) transform $r$ into the eigenbasis, $\hat r = Q^\top r$; (ii) scale each mode,
$\hat e_i = p(\lambda_i)\,\hat r_i$; (iii) transform back, $Pr = Q\hat e$. For a batch of
residuals $R\in\R^{N\times B}$ this is two dense matrix products. Every step is linear, so
the map $r\mapsto Pr$ is differentiable and gradients pass through it unchanged.

\paragraph{Cost.}
The eigendecomposition is a one-time $O(N^3)$ setup (performed in double precision; for the
grids used here $N$ is a few thousand, i.e.\ seconds), after which each application is
$O(N^2 B)$. Only $Q$ and the vector $p(\Lambda)$ need to be cached.

\paragraph{Diagnostics.}
The preconditioned spectrum is available in closed form from the cached eigenvalues,
\begin{equation}
  g(\lambda_i) = \lambda_i\,p(\lambda_i),
  \qquad
  \kk(PA) = \frac{\max_i g(\lambda_i)}{\min_i g(\lambda_i)},
  \qquad
  \kk(H_P) = \kk(PA)^2,
\end{equation}
which provides the exact conditioning of each configuration---the natural common axis when
comparing the two families or reporting how conditioning varies with the strength parameter.

\begin{remark}[Regular-grid shortcut]
On the regular grid the eigenvectors of $A$ are the tensor sine modes and $\Lambda$ is given
in closed form by \eqref{eq:eig2d}. Hence $Q^\top$ and $Q$ in \eqref{eq:apply} are
two-dimensional discrete sine transforms, and the dense eigendecomposition can be replaced by
an $O(N\log N)$ transform if larger grids are needed. At the resolutions considered here the
dense route is simplest and is used throughout.
\end{remark}

\end{document}
